% =============================================================================
%  RL-H4-01 | Draft pendahuluan, related work, dan metode
% -----------------------------------------------------------------------------
%  Role    : RL — Research Lead
%  Sprint  : H4 | PARALLEL (sisa waktu di luar RL-H4-02)
%  Ticket  : (isi dari papan kerja)          Status: STARTER ONLY
%  Target  : reports/sprint1/01_pendahuluan.tex, 02_related_work.tex, 03_metode.tex
%            (boleh satu file dulu; pisah sebelum H7)
%  Input   : brief v2 §1–§6, docs/positioning.pdf, docs/lit/
%  Output  : draft 3 bagian
%  Handoff : RL-H7-01
%  DoD     : [ ] semua klaim novelty memakai "sepengetahuan kami"
%            [ ] tidak ada klaim kategori JANGAN DIKLAIM
%            [ ] persamaan Modul 1 sama dengan brief §5
% =============================================================================
\documentclass[11pt,a4paper]{article}
% =============================================================================
%  00_preamble.tex | shared preamble untuk semua deliverable RL
%  Job     : bukan job card; dipakai oleh RL-H1-01 .. RL-H7-01
%  Target  : docs/ (salin bersama file .tex)
%  Catatan : kompilasi dengan pdflatex. Untuk XeLaTeX/LuaLaTeX, hapus inputenc.
% =============================================================================

\usepackage[utf8]{inputenc}
\usepackage[T1]{fontenc}
\usepackage{lmodern}
\usepackage[a4paper,margin=2.4cm]{geometry}
\usepackage{amsmath,amssymb}
\usepackage{booktabs}
\usepackage{longtable}
\usepackage{array}
\usepackage{graphicx}
\usepackage{xcolor}
\usepackage{enumitem}
\usepackage{float}
\usepackage{caption}
\usepackage{url}
\usepackage{microtype}
\usepackage[colorlinks=true,linkcolor=blue!50!black,urlcolor=blue!50!black,citecolor=blue!50!black]{hyperref}

% --- helper wajib -----------------------------------------------------------
% Path/nama file aman (underscore, garis miring) di dalam teks.
\newcommand{\fpath}[1]{\texttt{\detokenize{#1}}}
% TODO yang mencolok. Jangan dihapus sebelum job selesai.
\newcommand{\todo}[1]{\textcolor{red!75!black}{\textbf{[TODO: #1]}}}
\newcommand{\jobid}[1]{\textsf{\textbf{#1}}}
\newcommand{\runid}[1]{\texttt{\detokenize{#1}}}
\newcommand{\claim}[1]{\textsf{#1}}
\newcommand{\good}{\textcolor{green!45!black}{$\checkmark$}}
\newcommand{\bad}{\textcolor{red!75!black}{$\times$}}
\newcommand{\maybe}{\textcolor{orange!80!black}{$\sim$}}

\newcolumntype{L}[1]{>{\raggedright\arraybackslash}p{#1}}
\newcolumntype{C}[1]{>{\centering\arraybackslash}p{#1}}

\setlength{\parindent}{0pt}
\setlength{\parskip}{4pt}
\captionsetup{font=small,labelfont=bf}
\setlist[itemize]{leftmargin=1.2em,itemsep=2pt,topsep=2pt}
\setlist[enumerate]{leftmargin=1.5em,itemsep=2pt,topsep=2pt}

% Judul job yang konsisten di semua dokumen.
\newcommand{\jobheader}[4]{%
  \begin{center}
  {\Large\bfseries #1}\\[2pt]
  {\large #2}\\[4pt]
  \begin{tabular}{@{}ll@{}}
  \toprule
  Role & #3 \\
  Sprint & #4 \\
  \bottomrule
  \end{tabular}
  \end{center}
}


\title{\jobid{RL-H4-01} | Draft Laporan Sprint 1\\
\large Pendahuluan, Related Work, Metode}
\author{RL \\ \small RENARO Lab}
\date{\todo{tanggal H4} \quad | \quad \textcolor{red!75!black}{DRAFT}}

\begin{document}
\maketitle

\section{Pendahuluan}
\todo{Tulis 4--6 paragraf:
(1) harga pangan antardaerah saling berkaitan, dependensi tidak tetap/serentak;
(2) iklim di daerah produksi memberi sinyal dengan jeda yang berubah;
(3) gap: regime laten vs iklim teramati, lag per edge, evaluasi regime;
(4) RQ0--RQ3;
(5) kontribusi berurutan benchmark $\to$ bukti empiris $\to$ metode;
(6) scope A + batas: bukan identifikasi kausal.}

\paragraph{Kalimat kontribusi yang diizinkan.}
\begin{itemize}
  \item ``Sepengetahuan kami, belum ada ST-GNN antarprovinsi untuk harga cabai
        Indonesia yang mengondisikan distribusi jeda per edge pada iklim teramati
        di wilayah asal.'' \todo{verifikasi}
  \item ``Kami menyediakan benchmark evaluasi terstratifikasi regime dengan fitur
        availability-safe, worst-regime error, dan uji misspesifikasi graph prior.''
  \item ``Kami membandingkan secara head-to-head pengondisian iklim teramati
        melawan regime laten berbasis harga.''
\end{itemize}

\section{Related work}
\todo{Ringkas 4 kelompok dari brief \S2 dan \fpath{docs/positioning.tex}:
(1) graph learning spatio-temporal; (2) lag/delay-aware graph;
(3) OOD/robustness spatio-temporal; (4) harga pangan dan ENSO di Indonesia.
Akhiri dengan 3 kalimat posisi; jangan mengklaim hal yang sudah ada.}

\section{Metode}
\subsection{Representasi data dan aturan ketersediaan}
Node = provinsi; fitur = harga historis, kalender, hujan terbobot produksi
(wilayah sendiri dan asal), RONI vintage. Setiap fitur membawa \texttt{available\_at};
hanya nilai dengan \texttt{available\_at} $\le t$ yang dipakai pada decision date $t$.
\todo{tautkan ke kontrak data dan availability.csv}

\subsection{Modul 1: climate-conditioned lag graph}
Untuk tujuan $i$, asal $j$, dan bin jeda $\ell$:
\begin{equation}
A^{(\ell)}_{ij,t}
= \operatorname{softmax}_{j,\ell}
\left[
\phi_\theta\left(h_{i,\le t},\, h_{j,\le t},\, c_{j,\le t},\, s_{ij},\, \ell\right)
\right],
\label{eq:module1}
\end{equation}
dengan $h_{i,\le t}$ state encoder temporal, $c_{j,\le t}$ kondisi iklim asal yang
tersedia pada $t$, $s_{ij}$ fitur edge statis (jarak, mask BPS, pangsa), dan
$\ell=1,\dots,L$ bin pentad. Agregasi:
\begin{equation}
m_i = \sum_{j,\ell} A^{(\ell)}_{ij,t}\, W_v\, z_j^{(\ell)}.
\end{equation}
\todo{jelaskan encoder GRU, encoder iklim, representasi berjeda, opsi null,
candidate mask, dan head. Sebut kompleksitas untuk 34 node.}

\subsection{Modul 2--3: stable/adaptive dan reliability gate}
\begin{equation}
z_t = g_t \odot H^{\mathrm{adapt}}_t + (1-g_t)\odot H^{\mathrm{stable}}_t,
\qquad
g_{t,i} = \sigma\left(\mathrm{MLP}([\text{shock}_i,\text{rel}_i,\text{shift}_i])\right).
\end{equation}
\todo{definisikan shock/rel/shift dan pastikan semuanya tersedia pada $t$.}

\subsection{Modul 4 (stretch)}
\todo{penalti sensitivitas perturbasi, risiko lintas-regime dari fitur,
regularisasi prior lunak. Tegaskan bukan kontribusi utama.}

\subsection{Baseline dan protokol}
\todo{daftar baseline wajib + ringkasan split, metrik, DM, regime, holdout.
Tautkan ke \fpath{docs/eval_protocol_v1.pdf}.}

\section{Batas klaim}
\begin{table}[H]
\centering
\small
\begin{tabular}{@{}L{2.4cm}L{11.8cm}@{}}
\toprule
Kategori & Rumusan \\
\midrule
AMAN & \todo{salin dari brief \S6 + positioning} \\
BERISIKO & \todo{beri kualifikasi} \\
JANGAN & graph adaptif/gating/invariance sebagai novelty; edge = arus/kausalitas;
``pertama memakai iklim pada GNN harga pertanian'' \\
\bottomrule
\end{tabular}
\end{table}

\end{document}
